\documentclass[10pt,a4paper]{amsart}
\usepackage[margin=19mm]{geometry}
\usepackage{amsmath,amssymb,mathtools}
\usepackage[scr=rsfs,cal=euler]{mathalfa}
\usepackage{xfrac}
\usepackage[unicode,hidelinks,bookmarks=false]{hyperref}
\newcommand{\ie}{i.e.\ }
\newcommand{\cf}{cf.\ }
\newcommand{\resp}{respectively\ }
\newcommand{\Subsection}{\subsection}
\providecommand{\nobreakdash}{\nobreak\hbox{-}}
\setlength{\emergencystretch}{2em}
\multlinegap0pt
\newtheorem{lemma}{Lemma}
\title{Random Stability of Random Variables}
\author{Andrey Sarantsev}
\date{Source: arXiv:2603.28093v4 (CC BY 4.0)}
\begin{document}
\maketitle

\noindent\emph{Source and licence.} Random Stability of Random Variables. Version arXiv:2603.28093v4 (CC BY 4.0), distributed by arXiv under the Creative Commons Attribution 4.0 International licence, http://creativecommons.org/licenses/by/4.0/, as recorded in the arXiv licence metadata for this exact version and shown on the abstract page https://arxiv.org/abs/2603.28093v4. Full licence notice: \texttt{licenses/arxiv-2603.28093-CC-BY-4.0.txt}.\par\smallskip

\noindent\textit{Editorial scope.} Counted unit: the article's lemma eq:deriv (source0.tex lines 732--738). The article's complete original proof of it is delivered with it (source0.tex lines 740--764, one \texttt{proof} environment of 25 lines). No mathematics is written, repaired or re-derived: the statement and the proof are the article's own lines, in the article's own notation, and no retained line is substituted or re-flowed.  No further source-local context is needed: every cross-reference of every retained line resolves inside the two delivered spans.  Nothing was omitted from any retained span, so the package carries no omission record.  No retained line contains a citation, so the wrapper carries no bibliography and no bibliography entry.  No retained line contains a figure environment, an image-inclusion command or drawing code, so no image is delivered and the package declares no asset dependency.  Editorial formatting only, in addition: an AMS \texttt{amsart} preamble that restates the article's own declarations and macros used by the retained lines, namely the environments \texttt{lemma} on separate counters, together with the standard packages those lines need.  The retained environments print in this document as Lemma 1 in order of appearance, whereas the article prints the same items with the numbering of its own full document.  The complete native source of the article as distributed in the arXiv e-print for this exact version is bundled unchanged as \texttt{sources/197409/source0.tex}.
\begin{lemma} Assume $\mathbb E[X] < \infty$. We can define $L$ as a complex-valued function on the half-plane $\mathbb H = \{z \in \mathbb C\mid \mathrm{Re}(z) > 0\}$. This function is analytic on $\mathbb H$, and we have: 
\begin{equation}
\label{eq:deriv}
\frac{1 - L(z)}{z} \to \mathbb E[X],\quad z \to 0,\quad \mathrm{Re}(z) > 0.
\end{equation}
\label{lemma:tech}
\end{lemma}

\begin{proof} For $z = u + \mathrm{i}v \in \mathbb H$, we write $|e^{-zX}| = e^{-uX}$. Therefore, $\mathbb E|e^{-zX}| = \mathbb E[e^{-uX}] < \infty$, and $\mathbb E[e^{-zX}]$ is well defined. It is complex analytic on $\mathbb H$. To this end, we need to prove in the complex analytic sense: $L'(z) = \mathbb E[-Xe^{-zX}]$. But this, in turn, can be proved as follows. Take two points $z, z' \in \mathbb H$ and draw a segment between them. It is parametrized as $w_t = zt + z'(1-t)$. We need to show 
\begin{equation}
\label{eq:FTC}
L(w_t) - L(z) = \int_{[z, w_t]}\mathbb E\left[-Xe^{-uX}\right]\,\mathrm{d}u.
\end{equation}
This integral over the segment is understood in the complex analytic sense. This can be written using real-valued integration as
$$
\int_0^t\mathbb E\left[-Xe^{-(zs + (1-s)z')X}\right](z' - z)\,\mathrm{d}s.
$$
Remove expectations for a moment. From complex analysis, we get:
$$
e^{-(zt + z'(1-t))X} - e^{-ztX} = \int_0^t\left(-Xe^{-w_sX}\right)(z' - z)\,\mathrm{d}s.
$$
We need only to show that we need an interchange of integration and expectation. This requires us to use the Fubini theorem. Usually, this theorem is stated for real-vaued functions, but it works equally well for complex-valued functions. We need:
\begin{equation}
\label{eq:Fubini}
\mathbb E\int_0^t\left|-Xe^{-w_sX}\right|(z'-z)\,\mathrm{d}s < \infty.
\end{equation}
Assuming $\mathrm{Re}(z) \le \mathrm{Re}(z')$ without loss of generality, 
$\mathrm{Re}(z) \le \mathrm{Re}(w_s)$ for all $s \in [0, 1]$. Hence
$$
\mathbb E|-Xe^{-w_sX}| = \mathbb E|Xe^{-\mathrm{Re}(z)X}| < \infty.
$$
This proves~\eqref{eq:Fubini}, and with it proves~\eqref{eq:FTC}. Thus $L$ is analytic on $\mathbb H$. The property~\eqref{eq:deriv} can be shown similarly, with $z = 0$, since $\mathbb E[Xe^{-zX}] = \mathbb E[X]$ for $z = 0$. 
\end{proof}

\end{document}
